
\documentclass{article}
\usepackage{colortbl}
\usepackage{makecell}
\usepackage{multirow}
\usepackage{supertabular}

\begin{document}

\newcounter{utterance}

\centering \large Interaction Transcript for game `clean\_up', experiment `0\_easy\_7obj\_lv', episode 2 with nemotron{-}3{-}ultra.
\vspace{24pt}

{ \footnotesize  \setcounter{utterance}{1}
\setlength{\tabcolsep}{0pt}
\begin{supertabular}{c@{$\;$}|p{.15\linewidth}@{}p{.15\linewidth}p{.15\linewidth}p{.15\linewidth}p{.15\linewidth}p{.15\linewidth}}
    \# & $\;$A & \multicolumn{4}{c}{Game Master} & $\;\:$B\\
    \hline

    \theutterance \stepcounter{utterance}  
    & & \multicolumn{4}{p{0.6\linewidth}}{
        \cellcolor[rgb]{0.9,0.9,0.9}{
            \makecell[{{p{\linewidth}}}]{
                \texttt{\tiny{[P1$\langle$GM]}}
                \texttt{Es esmu jūsu spēles vadītājs, un jūs spēlējat sadarbības spēli ar šādu režģi kā spēles laukumu:} \\
\\ 
\texttt{```} \\
\texttt{ 1234567} \\
\texttt{╔═══════╗} \\
\texttt{║◌T◌◌◌◌◌║ 1} \\
\texttt{║A◌◌◌◌◌D║ 2} \\
\texttt{║◌◌◌◌◌┌─╢ 3} \\
\texttt{╟──┐O◌│◌║ 4} \\
\texttt{║◌P├──┼─╢ 5} \\
\texttt{║◌S│M◌│◌║ 6} \\
\texttt{║◌◌│◌◌│◌║ 7} \\
\texttt{╚══╧══╧═╝} \\
\\ 
\texttt{```} \\
\\ 
\texttt{* Augšējā malā ir attēlotas x koordinātas, kas pieaug pa labi, un labajā malā y koordinātas, kas pieaug uz leju.} \\
\texttt{* Uz jūsu režģa ir nejauši izvietoti šādi objekti: 'P', 'O', 'T', 'S', 'D', 'A', 'M'.} \\
\\ 
\\ 
\texttt{Otrs spēlētājs redz spēles laukuma variāciju, kurā objekti ir izvietoti citās nejaušās vietās. Jūs nevarat redzēt otra spēlētāja laukumu, un viņš nevar redzēt jūsējo.} \\
\\ 
\\ 
\texttt{**Mērķis:**} \\
\\ 
\texttt{Abiem spēlētājiem ir jāpārvieto objekti uz sava fona tā, lai vienādi objekti nonāktu vienādās koordinātēs. Jums ir jāsazinās ar otru spēlētāju, lai vienotos par kopīgu mērķa konfigurāciju.} \\
\\ 
\\ 
\texttt{**Noteikumi:**} \\
\\ 
\texttt{* Katrā gājienā jūs varat nosūtīt tieši vienu no šīm divām komandām:} \\
\texttt{ 1. `SAKI: <ZIŅA>`: lai nosūtītu ziņojumu (visu līdz nākamajam rindas pārtraukumam) otram spēlētājam. Es to pārsūtīšu jūsu partnerim.} \\
\texttt{2. `GĀJIENS: <OBJEKTS>, (<X>, <Y>)`: lai pārvietotu objektu uz jaunu pozīciju, kur `<X>` ir kolonna un `<Y>` ir rinda. Es jūs informēšu, vai jūsu gājiens bija veiksmīgs vai nē.} \\
\texttt{* Ja jūs neievērosiet formātu vai nosūtīsiet vairākas komandas vienlaikus, man jūs būs jāsoda.} \\
\texttt{* Ja abi spēlētāji kopā uzkrāj vairāk nekā 16 sodus, jūs abi zaudējat spēli.} \\
\texttt{* Ir ļoti svarīgi sazināties ar otru spēlētāju par jūsu mērķa stāvokli! *Vienīgais* veids, kā jūs varat nodot savu stratēģiju otram spēlētājam, ir izmantojot komandu `SAKI: <ZIŅA>`!} \\
\\ 
\\ 
\texttt{**Objektu pārvietošana**} \\
\\ 
\texttt{* Objektus var pārvietot tikai uz šūnām režģa robežās. Mērķa šūnai jābūt tukšai, t.i., tajā jābūt tikai simbolam '◌'.} \\
\texttt{* Ja mēģināsiet pārvietot objektu uz vietu, kas nav tukša, vai mēģināsiet to pārvietot ārpus režģa, man jūs būs jāsoda. Jūs saņemat vēl vienu mēģinājumu.} \\
\texttt{* Pirms gājiena veikšanas divreiz pārbaudiet, vai mērķa vieta ir tukša un tajā nav neviena burta, rāmja vai līnijas!} \\
\\ 
\\ 
\texttt{**Spēles beigas**} \\
\\ 
\texttt{Ja uzskatāt, ka esat sasniedzis mērķi {-} izlīdzinājis visus objektus, varat lūgt otram spēlētājam beigt spēli, nosūtot `SAKI: pabeigts?`. Ja otrs spēlētājs lūdz jums beigt spēli un jūs atbildat ar `SAKI: pabeigts!`, spēle beigsies.} \\
\\ 
\texttt{Abi spēlētāji uzvar, ja spēle beidzas 28 raundos, kur viens raunds ir definēts kā divi spēlētāji, kas katrs nosūta derīgu komandu.} \\
\\ 
\texttt{**Punktu skaitīšana:**} \\
\\ 
\texttt{Jo tuvāk vienādi objekti atrodas abos spēles laukumos, jo vairāk punktu jūs saņemat. Sodi samazina jūsu punktus. Vai varat pārspēt rekordu?} \\
\\ 
\\ 
\texttt{Lūdzu, nosūtiet ziņu otram spēlētājam, lai sāktu spēli!} \\
            }
        }
    }
    & \\ \\

    \theutterance \stepcounter{utterance}  
    & \multicolumn{4}{p{0.6\linewidth}}{
        \cellcolor[rgb]{0.9,0.9,0.9}{
            \makecell[{{p{\linewidth}}}]{
                \texttt{\tiny{[P1$\rangle$GM]}}
                \texttt{SAKI: Sveicināti! Sākam koordināciju mainīšanu. Man šajā brīdī objekti atrodas šādās pozīcijās:} \\
\texttt{T: (2, 1)} \\
\texttt{A: (1, 2)} \\
\texttt{D: (7, 2)} \\
\texttt{O: (5, 4)} \\
\texttt{P: (2, 5)} \\
\texttt{S: (2, 6)} \\
\texttt{M: (4, 6)} \\
\\ 
\texttt{Kurp jūsu objekti atrodas? Varbūt varam vienoties, lai mērķis būtu grupēt visu pa kreiso pusē (piemēram, kolonnās 1{-}3), lai vieglāk sinhronizētu?} \\
            }
        }
    }
    & & \\ \\

    \theutterance \stepcounter{utterance}  
    & & & \multicolumn{2}{p{0.3\linewidth}}{
        \cellcolor[rgb]{0.9,0.9,0.9}{
            \makecell[{{p{\linewidth}}}]{
                \texttt{\tiny{[GM$|$GM]}}
                \texttt{SAKI: Sveicināti! Sākam koordināciju mainīšanu. Man šajā brīdī objekti atrodas šādās pozīcijās:} \\
\texttt{T: (2, 1)} \\
\texttt{A: (1, 2)} \\
\texttt{D: (7, 2)} \\
\texttt{O: (5, 4)} \\
\texttt{P: (2, 5)} \\
\texttt{S: (2, 6)} \\
\texttt{M: (4, 6)} \\
\\ 
\texttt{Kurp jūsu objekti atrodas? Varbūt varam vienoties, lai mērķis būtu grupēt visu pa kreiso pusē (piemēram, kolonnās 1{-}3), lai vieglāk sinhronizētu?} \\
            }
        }
    }
    & & \\ \\

    \theutterance \stepcounter{utterance}  
    & & & \multicolumn{2}{p{0.3\linewidth}}{
        \cellcolor[rgb]{0.9,0.9,0.9}{
            \makecell[{{p{\linewidth}}}]{
                \texttt{\tiny{[GM$|$GM]}}
                \texttt{Invalid format: tail is not empty: '} \\
\texttt{T: (2, 1)} \\
\texttt{A: (1, 2)} \\
\texttt{D: (7, 2)} \\
\texttt{O: (5, 4)} \\
\texttt{P: (2, 5)} \\
\texttt{S: (2, 6)} \\
\texttt{M: (4, 6)} \\
\\ 
\texttt{Kurp jūsu objekti atrodas? Varbūt varam vienoties, lai mērķis būtu grupēt visu pa kreiso pusē (piemēram, kolonnās 1{-}3), lai vieglāk sinhronizētu?} \\
            }
        }
    }
    & & \\ \\

    \theutterance \stepcounter{utterance}  
    & & \multicolumn{4}{p{0.6\linewidth}}{
        \cellcolor[rgb]{0.9,0.9,0.9}{
            \makecell[{{p{\linewidth}}}]{
                \texttt{\tiny{[P1$\langle$GM]}}
                \texttt{Sods: Jūsu ziņojumā nedrīkst būt nekas pēc komandas!} \\
\\ 
\texttt{Pārliecinieties, ka jūsu atbildē ir tikai `SAKI: <ZIŅA>` vai `GĀJIENS: <OBJEKTS>, (<X>, <Y>)`, un nekas cits!} \\
\\ 
\texttt{Jūs kopā esat uzkrājuši 1 no 16 sodiem.} \\
\texttt{Lūdzu, mēģiniet vēlreiz!} \\
            }
        }
    }
    & \\ \\

    \theutterance \stepcounter{utterance}  
    & \multicolumn{4}{p{0.6\linewidth}}{
        \cellcolor[rgb]{0.9,0.9,0.9}{
            \makecell[{{p{\linewidth}}}]{
                \texttt{\tiny{[P1$\rangle$GM]}}
                \texttt{SAKI: Sveicināti! Mani objekti: T(2,1), A(1,2), D(7,2), O(5,4), P(2,5), S(2,6), M(4,6). Kur atrodas jūsu objekti? Varbūt vienojamies par kopīgu mērķa izkārtojumu?} \\
            }
        }
    }
    & & \\ \\

    \theutterance \stepcounter{utterance}  
    & & & \multicolumn{2}{p{0.3\linewidth}}{
        \cellcolor[rgb]{0.9,0.9,0.9}{
            \makecell[{{p{\linewidth}}}]{
                \texttt{\tiny{[GM$|$GM]}}
                \texttt{SAKI: Sveicināti! Mani objekti: T(2,1), A(1,2), D(7,2), O(5,4), P(2,5), S(2,6), M(4,6). Kur atrodas jūsu objekti? Varbūt vienojamies par kopīgu mērķa izkārtojumu?} \\
            }
        }
    }
    & & \\ \\

    \theutterance \stepcounter{utterance}  
    & & \multicolumn{4}{p{0.6\linewidth}}{
        \cellcolor[rgb]{0.9,0.9,0.9}{
            \makecell[{{p{\linewidth}}}]{
                \texttt{\tiny{[GM$\rangle$P2]}}
                \texttt{Es esmu jūsu spēles vadītājs, un jūs spēlējat sadarbības spēli ar šādu režģi kā spēles laukumu:} \\
\\ 
\texttt{```} \\
\texttt{ 1234567} \\
\texttt{╔═══════╗} \\
\texttt{║◌◌P◌A◌O║ 1} \\
\texttt{║S◌◌◌M◌◌║ 2} \\
\texttt{║◌◌◌◌D┌─╢ 3} \\
\texttt{╟──┐◌◌│◌║ 4} \\
\texttt{║◌◌├──┼─╢ 5} \\
\texttt{║◌◌│◌◌│T║ 6} \\
\texttt{║◌◌│◌◌│◌║ 7} \\
\texttt{╚══╧══╧═╝} \\
\\ 
\texttt{```} \\
\\ 
\texttt{* Augšējā malā ir attēlotas x koordinātas, kas pieaug pa labi, un labajā malā y koordinātas, kas pieaug uz leju.} \\
\texttt{* Uz jūsu režģa ir nejauši izvietoti šādi objekti: 'P', 'O', 'T', 'S', 'D', 'A', 'M'.} \\
\\ 
\\ 
\texttt{Otrs spēlētājs redz spēles laukuma variāciju, kurā objekti ir izvietoti citās nejaušās vietās. Jūs nevarat redzēt otra spēlētāja laukumu, un viņš nevar redzēt jūsējo.} \\
\\ 
\\ 
\texttt{**Mērķis:**} \\
\\ 
\texttt{Abiem spēlētājiem ir jāpārvieto objekti uz sava fona tā, lai vienādi objekti nonāktu vienādās koordinātēs. Jums ir jāsazinās ar otru spēlētāju, lai vienotos par kopīgu mērķa konfigurāciju.} \\
\\ 
\\ 
\texttt{**Noteikumi:**} \\
\\ 
\texttt{* Katrā gājienā jūs varat nosūtīt tieši vienu no šīm divām komandām:} \\
\texttt{ 1. `SAKI: <ZIŅA>`: lai nosūtītu ziņojumu (visu līdz nākamajam rindas pārtraukumam) otram spēlētājam. Es to pārsūtīšu jūsu partnerim.} \\
\texttt{2. `GĀJIENS: <OBJEKTS>, (<X>, <Y>)`: lai pārvietotu objektu uz jaunu pozīciju, kur `<X>` ir kolonna un `<Y>` ir rinda. Es jūs informēšu, vai jūsu gājiens bija veiksmīgs vai nē.} \\
\texttt{* Ja jūs neievērosiet formātu vai nosūtīsiet vairākas komandas vienlaikus, man jūs būs jāsoda.} \\
\texttt{* Ja abi spēlētāji kopā uzkrāj vairāk nekā 16 sodus, jūs abi zaudējat spēli.} \\
\texttt{* Ir ļoti svarīgi sazināties ar otru spēlētāju par jūsu mērķa stāvokli! *Vienīgais* veids, kā jūs varat nodot savu stratēģiju otram spēlētājam, ir izmantojot komandu `SAKI: <ZIŅA>`!} \\
\\ 
\\ 
\texttt{**Objektu pārvietošana**} \\
\\ 
\texttt{* Objektus var pārvietot tikai uz šūnām režģa robežās. Mērķa šūnai jābūt tukšai, t.i., tajā jābūt tikai simbolam '◌'.} \\
\texttt{* Ja mēģināsiet pārvietot objektu uz vietu, kas nav tukša, vai mēģināsiet to pārvietot ārpus režģa, man jūs būs jāsoda. Jūs saņemat vēl vienu mēģinājumu.} \\
\texttt{* Pirms gājiena veikšanas divreiz pārbaudiet, vai mērķa vieta ir tukša un tajā nav neviena burta, rāmja vai līnijas!} \\
\\ 
\\ 
\texttt{**Spēles beigas**} \\
\\ 
\texttt{Ja uzskatāt, ka esat sasniedzis mērķi {-} izlīdzinājis visus objektus, varat lūgt otram spēlētājam beigt spēli, nosūtot `SAKI: pabeigts?`. Ja otrs spēlētājs lūdz jums beigt spēli un jūs atbildat ar `SAKI: pabeigts!`, spēle beigsies.} \\
\\ 
\texttt{Abi spēlētāji uzvar, ja spēle beidzas 28 raundos, kur viens raunds ir definēts kā divi spēlētāji, kas katrs nosūta derīgu komandu.} \\
\\ 
\texttt{**Punktu skaitīšana:**} \\
\\ 
\texttt{Jo tuvāk vienādi objekti atrodas abos spēles laukumos, jo vairāk punktu jūs saņemat. Sodi samazina jūsu punktus. Vai varat pārspēt rekordu?} \\
\\ 
\texttt{Otrs spēlētājs sāka spēli, nosūtot šo ziņojumu:} \\
\\ 
\texttt{"Sveicināti! Mani objekti: T(2,1), A(1,2), D(7,2), O(5,4), P(2,5), S(2,6), M(4,6). Kur atrodas jūsu objekti? Varbūt vienojamies par kopīgu mērķa izkārtojumu?"} \\
\\ 
\texttt{Kāda ir jūsu pirmā komanda?} \\
            }
        }
    }
    & \\ \\

    \theutterance \stepcounter{utterance}  
    & & & \multicolumn{4}{p{0.6\linewidth}}{
        \cellcolor[rgb]{0.9,0.9,0.9}{
            \makecell[{{p{\linewidth}}}]{
                \texttt{\tiny{[GM$\langle$P2]}}
                \texttt{SAKI: Sveicināti! Mani objekti: P(3,1), A(5,1), O(7,1), S(1,2), M(5,2), D(4,3), T(6,6). Man šķiet, ka mums ir jāsaskaņojas par kopīgu mērķa izkārtojumu, lai abiem laukumiem objektu koordinātes sakrītu. Jums kādi varianti patīk? Varbūt izvietojam tos rindās 1{-}3, lai izvairītos no sienu?} \\
            }
        }
    }
     \\ \\

    \theutterance \stepcounter{utterance}  
    & & & \multicolumn{2}{p{0.3\linewidth}}{
        \cellcolor[rgb]{0.9,0.9,0.9}{
            \makecell[{{p{\linewidth}}}]{
                \texttt{\tiny{[GM$|$GM]}}
                \texttt{SAKI: Sveicināti! Mani objekti: P(3,1), A(5,1), O(7,1), S(1,2), M(5,2), D(4,3), T(6,6). Man šķiet, ka mums ir jāsaskaņojas par kopīgu mērķa izkārtojumu, lai abiem laukumiem objektu koordinātes sakrītu. Jums kādi varianti patīk? Varbūt izvietojam tos rindās 1{-}3, lai izvairītos no sienu?} \\
            }
        }
    }
    & & \\ \\

    \theutterance \stepcounter{utterance}  
    & & \multicolumn{4}{p{0.6\linewidth}}{
        \cellcolor[rgb]{0.9,0.9,0.9}{
            \makecell[{{p{\linewidth}}}]{
                \texttt{\tiny{[P1$\langle$GM]}}
                \texttt{Jūsu ziņojums ir nodots otram spēlētājam.} \\
\\ 
\\ 
\texttt{Pašlaik jūs spēlējat kārtu 1 no maksimāli 28.} \\
\texttt{Jūs kopā esat uzkrājuši 1 no 16 sodiem.} \\
\texttt{Otrs spēlētājs nosūtīja šo ziņojumu:} \\
\\ 
\\ 
\texttt{"Sveicināti! Mani objekti: P(3,1), A(5,1), O(7,1), S(1,2), M(5,2), D(4,3), T(6,6). Man šķiet, ka mums ir jāsaskaņojas par kopīgu mērķa izkārtojumu, lai abiem laukumiem objektu koordinātes sakrītu. Jums kādi varianti patīk? Varbūt izvietojam tos rindās 1{-}3, lai izvairītos no sienu?"} \\
\\ 
\\ 
\texttt{Kāda ir jūsu nākamā komanda?} \\
            }
        }
    }
    & \\ \\

    \theutterance \stepcounter{utterance}  
    & \multicolumn{4}{p{0.6\linewidth}}{
        \cellcolor[rgb]{0.9,0.9,0.9}{
            \makecell[{{p{\linewidth}}}]{
                \texttt{\tiny{[P1$\rangle$GM]}}
                \texttt{SAKI: Labi, rindās 1{-}3 ir brīvs lauks (mans laukums 1{-}3 rindās tikai kolonnās 1,2,5,7). Piemērojam secību pa kreiso augšējo stūri: Rinda 1: A(1,1), D(2,1), M(3,1), O(4,1), P(5,1), S(6,1), T(7,1). Rindas 2{-}3 paliek tukšas rezervē. Vai šis variants pietiek? Ja jau, sākam pārvietot. Mans pirmais gājiens: A uz (1,1).} \\
            }
        }
    }
    & & \\ \\

    \theutterance \stepcounter{utterance}  
    & & & \multicolumn{2}{p{0.3\linewidth}}{
        \cellcolor[rgb]{0.9,0.9,0.9}{
            \makecell[{{p{\linewidth}}}]{
                \texttt{\tiny{[GM$|$GM]}}
                \texttt{SAKI: Labi, rindās 1{-}3 ir brīvs lauks (mans laukums 1{-}3 rindās tikai kolonnās 1,2,5,7). Piemērojam secību pa kreiso augšējo stūri: Rinda 1: A(1,1), D(2,1), M(3,1), O(4,1), P(5,1), S(6,1), T(7,1). Rindas 2{-}3 paliek tukšas rezervē. Vai šis variants pietiek? Ja jau, sākam pārvietot. Mans pirmais gājiens: A uz (1,1).} \\
            }
        }
    }
    & & \\ \\

    \theutterance \stepcounter{utterance}  
    & & \multicolumn{4}{p{0.6\linewidth}}{
        \cellcolor[rgb]{0.9,0.9,0.9}{
            \makecell[{{p{\linewidth}}}]{
                \texttt{\tiny{[GM$\rangle$P2]}}
                \texttt{Jūsu ziņojums ir nodots otram spēlētājam.} \\
\\ 
\\ 
\texttt{Pašlaik jūs spēlējat kārtu 2 no maksimāli 28.} \\
\texttt{Jūs kopā esat uzkrājuši 1 no 16 sodiem.} \\
\texttt{Otrs spēlētājs nosūtīja šo ziņojumu:} \\
\\ 
\\ 
\texttt{"Labi, rindās 1{-}3 ir brīvs lauks (mans laukums 1{-}3 rindās tikai kolonnās 1,2,5,7). Piemērojam secību pa kreiso augšējo stūri: Rinda 1: A(1,1), D(2,1), M(3,1), O(4,1), P(5,1), S(6,1), T(7,1). Rindas 2{-}3 paliek tukšas rezervē. Vai šis variants pietiek? Ja jau, sākam pārvietot. Mans pirmais gājiens: A uz (1,1)."} \\
\\ 
\\ 
\texttt{Kāda ir jūsu nākamā komanda?} \\
            }
        }
    }
    & \\ \\

    \theutterance \stepcounter{utterance}  
    & & & \multicolumn{4}{p{0.6\linewidth}}{
        \cellcolor[rgb]{0.9,0.9,0.9}{
            \makecell[{{p{\linewidth}}}]{
                \texttt{\tiny{[GM$\langle$P2]}}
                \texttt{GĀJIENS: A, (5, 1)} \\
            }
        }
    }
     \\ \\

    \theutterance \stepcounter{utterance}  
    & & & \multicolumn{2}{p{0.3\linewidth}}{
        \cellcolor[rgb]{0.9,0.9,0.9}{
            \makecell[{{p{\linewidth}}}]{
                \texttt{\tiny{[GM$|$GM]}}
                \texttt{GĀJIENS: A, (5, 1)} \\
            }
        }
    }
    & & \\ \\

    \theutterance \stepcounter{utterance}  
    & & & \multicolumn{2}{p{0.3\linewidth}}{
        \cellcolor[rgb]{0.9,0.9,0.9}{
            \makecell[{{p{\linewidth}}}]{
                \texttt{\tiny{[GM$|$GM]}}
                \texttt{Sods: (5,1) nav tukšs, bet satur 'A'.} \\
\\ 
\texttt{Jūs kopā esat uzkrājuši 2 no 16 sodiem.} \\
\texttt{Lūdzu, mēģiniet vēlreiz!} \\
            }
        }
    }
    & & \\ \\

    \theutterance \stepcounter{utterance}  
    & & \multicolumn{4}{p{0.6\linewidth}}{
        \cellcolor[rgb]{0.9,0.9,0.9}{
            \makecell[{{p{\linewidth}}}]{
                \texttt{\tiny{[GM$\rangle$P2]}}
                \texttt{Sods: (5,1) nav tukšs, bet satur 'A'.} \\
\\ 
\texttt{Jūs kopā esat uzkrājuši 2 no 16 sodiem.} \\
\texttt{Lūdzu, mēģiniet vēlreiz!} \\
            }
        }
    }
    & \\ \\

    \theutterance \stepcounter{utterance}  
    & & & \multicolumn{4}{p{0.6\linewidth}}{
        \cellcolor[rgb]{0.9,0.9,0.9}{
            \makecell[{{p{\linewidth}}}]{
                \texttt{\tiny{[GM$\langle$P2]}}
                \texttt{GĀJIENS: A, (1, 1)} \\
            }
        }
    }
     \\ \\

    \theutterance \stepcounter{utterance}  
    & & & \multicolumn{2}{p{0.3\linewidth}}{
        \cellcolor[rgb]{0.9,0.9,0.9}{
            \makecell[{{p{\linewidth}}}]{
                \texttt{\tiny{[GM$|$GM]}}
                \texttt{GĀJIENS: A, (1, 1)} \\
            }
        }
    }
    & & \\ \\

    \theutterance \stepcounter{utterance}  
    & & & \multicolumn{2}{p{0.3\linewidth}}{
        \cellcolor[rgb]{0.9,0.9,0.9}{
            \makecell[{{p{\linewidth}}}]{
                \texttt{\tiny{[GM$|$GM]}}
                \texttt{A moved to (1, 1)} \\
            }
        }
    }
    & & \\ \\

    \theutterance \stepcounter{utterance}  
    & & \multicolumn{4}{p{0.6\linewidth}}{
        \cellcolor[rgb]{0.9,0.9,0.9}{
            \makecell[{{p{\linewidth}}}]{
                \texttt{\tiny{[P1$\langle$GM]}}
                \texttt{Jūsu ziņojums ir nodots otram spēlētājam.} \\
\\ 
\\ 
\texttt{Pašlaik jūs spēlējat kārtu 2 no maksimāli 28.} \\
\texttt{Jūs kopā esat uzkrājuši 2 no 16 sodiem.} \\
\texttt{Otrs spēlētājs pārvietoja objektu savā režģī.} \\
\\ 
\\ 
\texttt{Kāda ir jūsu nākamā komanda?} \\
            }
        }
    }
    & \\ \\

    \theutterance \stepcounter{utterance}  
    & \multicolumn{4}{p{0.6\linewidth}}{
        \cellcolor[rgb]{0.9,0.9,0.9}{
            \makecell[{{p{\linewidth}}}]{
                \texttt{\tiny{[P1$\rangle$GM]}}
                \texttt{GĀJIENS: A, (1, 1)} \\
            }
        }
    }
    & & \\ \\

    \theutterance \stepcounter{utterance}  
    & & & \multicolumn{2}{p{0.3\linewidth}}{
        \cellcolor[rgb]{0.9,0.9,0.9}{
            \makecell[{{p{\linewidth}}}]{
                \texttt{\tiny{[GM$|$GM]}}
                \texttt{GĀJIENS: A, (1, 1)} \\
            }
        }
    }
    & & \\ \\

    \theutterance \stepcounter{utterance}  
    & & & \multicolumn{2}{p{0.3\linewidth}}{
        \cellcolor[rgb]{0.9,0.9,0.9}{
            \makecell[{{p{\linewidth}}}]{
                \texttt{\tiny{[GM$|$GM]}}
                \texttt{A moved to (1, 1)} \\
            }
        }
    }
    & & \\ \\

    \theutterance \stepcounter{utterance}  
    & & \multicolumn{4}{p{0.6\linewidth}}{
        \cellcolor[rgb]{0.9,0.9,0.9}{
            \makecell[{{p{\linewidth}}}]{
                \texttt{\tiny{[GM$\rangle$P2]}}
                \texttt{Veiksmīgi pārvietots 'A' uz (1,1). Jūsu atjauninātais režģis izskatās šādi:} \\
\texttt{```} \\
\texttt{ 1234567} \\
\texttt{╔═══════╗} \\
\texttt{║A◌P◌◌◌O║ 1} \\
\texttt{║S◌◌◌M◌◌║ 2} \\
\texttt{║◌◌◌◌D┌─╢ 3} \\
\texttt{╟──┐◌◌│◌║ 4} \\
\texttt{║◌◌├──┼─╢ 5} \\
\texttt{║◌◌│◌◌│T║ 6} \\
\texttt{║◌◌│◌◌│◌║ 7} \\
\texttt{╚══╧══╧═╝} \\
\texttt{```} \\
\texttt{Pašlaik jūs spēlējat kārtu 3 no maksimāli 28.} \\
\texttt{Jūs kopā esat uzkrājuši 2 no 16 sodiem.} \\
\texttt{Otrs spēlētājs pārvietoja objektu savā režģī.} \\
\\ 
\\ 
\texttt{Kāda ir jūsu nākamā komanda?} \\
            }
        }
    }
    & \\ \\

    \theutterance \stepcounter{utterance}  
    & & & \multicolumn{4}{p{0.6\linewidth}}{
        \cellcolor[rgb]{0.9,0.9,0.9}{
            \makecell[{{p{\linewidth}}}]{
                \texttt{\tiny{[GM$\langle$P2]}}
                \texttt{GĀJIENS: D, (2, 1)} \\
            }
        }
    }
     \\ \\

    \theutterance \stepcounter{utterance}  
    & & & \multicolumn{2}{p{0.3\linewidth}}{
        \cellcolor[rgb]{0.9,0.9,0.9}{
            \makecell[{{p{\linewidth}}}]{
                \texttt{\tiny{[GM$|$GM]}}
                \texttt{GĀJIENS: D, (2, 1)} \\
            }
        }
    }
    & & \\ \\

    \theutterance \stepcounter{utterance}  
    & & & \multicolumn{2}{p{0.3\linewidth}}{
        \cellcolor[rgb]{0.9,0.9,0.9}{
            \makecell[{{p{\linewidth}}}]{
                \texttt{\tiny{[GM$|$GM]}}
                \texttt{D moved to (2, 1)} \\
            }
        }
    }
    & & \\ \\

    \theutterance \stepcounter{utterance}  
    & & \multicolumn{4}{p{0.6\linewidth}}{
        \cellcolor[rgb]{0.9,0.9,0.9}{
            \makecell[{{p{\linewidth}}}]{
                \texttt{\tiny{[P1$\langle$GM]}}
                \texttt{Veiksmīgi pārvietots 'A' uz (1,1). Jūsu atjauninātais režģis izskatās šādi:} \\
\texttt{```} \\
\texttt{ 1234567} \\
\texttt{╔═══════╗} \\
\texttt{║AT◌◌◌◌◌║ 1} \\
\texttt{║◌◌◌◌◌◌D║ 2} \\
\texttt{║◌◌◌◌◌┌─╢ 3} \\
\texttt{╟──┐O◌│◌║ 4} \\
\texttt{║◌P├──┼─╢ 5} \\
\texttt{║◌S│M◌│◌║ 6} \\
\texttt{║◌◌│◌◌│◌║ 7} \\
\texttt{╚══╧══╧═╝} \\
\texttt{```} \\
\texttt{Pašlaik jūs spēlējat kārtu 3 no maksimāli 28.} \\
\texttt{Jūs kopā esat uzkrājuši 2 no 16 sodiem.} \\
\texttt{Otrs spēlētājs pārvietoja objektu savā režģī.} \\
\\ 
\\ 
\texttt{Kāda ir jūsu nākamā komanda?} \\
            }
        }
    }
    & \\ \\

    \theutterance \stepcounter{utterance}  
    & \multicolumn{4}{p{0.6\linewidth}}{
        \cellcolor[rgb]{0.9,0.9,0.9}{
            \makecell[{{p{\linewidth}}}]{
                \texttt{\tiny{[P1$\rangle$GM]}}
                \texttt{GĀJIENS: D, (2, 1)} \\
            }
        }
    }
    & & \\ \\

    \theutterance \stepcounter{utterance}  
    & & & \multicolumn{2}{p{0.3\linewidth}}{
        \cellcolor[rgb]{0.9,0.9,0.9}{
            \makecell[{{p{\linewidth}}}]{
                \texttt{\tiny{[GM$|$GM]}}
                \texttt{GĀJIENS: D, (2, 1)} \\
            }
        }
    }
    & & \\ \\

    \theutterance \stepcounter{utterance}  
    & & & \multicolumn{2}{p{0.3\linewidth}}{
        \cellcolor[rgb]{0.9,0.9,0.9}{
            \makecell[{{p{\linewidth}}}]{
                \texttt{\tiny{[GM$|$GM]}}
                \texttt{Sods: (2,1) nav tukšs, bet satur 'T'.} \\
\\ 
\texttt{Jūs kopā esat uzkrājuši 3 no 16 sodiem.} \\
\texttt{Lūdzu, mēģiniet vēlreiz!} \\
            }
        }
    }
    & & \\ \\

    \theutterance \stepcounter{utterance}  
    & & \multicolumn{4}{p{0.6\linewidth}}{
        \cellcolor[rgb]{0.9,0.9,0.9}{
            \makecell[{{p{\linewidth}}}]{
                \texttt{\tiny{[P1$\langle$GM]}}
                \texttt{Sods: (2,1) nav tukšs, bet satur 'T'.} \\
\\ 
\texttt{Jūs kopā esat uzkrājuši 3 no 16 sodiem.} \\
\texttt{Lūdzu, mēģiniet vēlreiz!} \\
            }
        }
    }
    & \\ \\

    \theutterance \stepcounter{utterance}  
    & \multicolumn{4}{p{0.6\linewidth}}{
        \cellcolor[rgb]{0.9,0.9,0.9}{
            \makecell[{{p{\linewidth}}}]{
                \texttt{\tiny{[P1$\rangle$GM]}}
                \texttt{SAKI: Atvainojos, kļūdījās ar D mērķi. T jau ir (2,1). Turpmāk pārvietojam T uz (3,1), pēc tam D uz (2,1). Mani nākamais gājiens: T uz (3,1).} \\
            }
        }
    }
    & & \\ \\

    \theutterance \stepcounter{utterance}  
    & & & \multicolumn{2}{p{0.3\linewidth}}{
        \cellcolor[rgb]{0.9,0.9,0.9}{
            \makecell[{{p{\linewidth}}}]{
                \texttt{\tiny{[GM$|$GM]}}
                \texttt{SAKI: Atvainojos, kļūdījās ar D mērķi. T jau ir (2,1). Turpmāk pārvietojam T uz (3,1), pēc tam D uz (2,1). Mani nākamais gājiens: T uz (3,1).} \\
            }
        }
    }
    & & \\ \\

    \theutterance \stepcounter{utterance}  
    & & \multicolumn{4}{p{0.6\linewidth}}{
        \cellcolor[rgb]{0.9,0.9,0.9}{
            \makecell[{{p{\linewidth}}}]{
                \texttt{\tiny{[GM$\rangle$P2]}}
                \texttt{Veiksmīgi pārvietots 'D' uz (2,1). Jūsu atjauninātais režģis izskatās šādi:} \\
\texttt{```} \\
\texttt{ 1234567} \\
\texttt{╔═══════╗} \\
\texttt{║ADP◌◌◌O║ 1} \\
\texttt{║S◌◌◌M◌◌║ 2} \\
\texttt{║◌◌◌◌◌┌─╢ 3} \\
\texttt{╟──┐◌◌│◌║ 4} \\
\texttt{║◌◌├──┼─╢ 5} \\
\texttt{║◌◌│◌◌│T║ 6} \\
\texttt{║◌◌│◌◌│◌║ 7} \\
\texttt{╚══╧══╧═╝} \\
\texttt{```} \\
\texttt{Pašlaik jūs spēlējat kārtu 4 no maksimāli 28.} \\
\texttt{Jūs kopā esat uzkrājuši 3 no 16 sodiem.} \\
\texttt{Otrs spēlētājs nosūtīja šo ziņojumu:} \\
\\ 
\\ 
\texttt{"Atvainojos, kļūdījās ar D mērķi. T jau ir (2,1). Turpmāk pārvietojam T uz (3,1), pēc tam D uz (2,1). Mani nākamais gājiens: T uz (3,1)."} \\
\\ 
\\ 
\texttt{Kāda ir jūsu nākamā komanda?} \\
            }
        }
    }
    & \\ \\

    \theutterance \stepcounter{utterance}  
    & & & \multicolumn{4}{p{0.6\linewidth}}{
        \cellcolor[rgb]{0.9,0.9,0.9}{
            \makecell[{{p{\linewidth}}}]{
                \texttt{\tiny{[GM$\langle$P2]}}
                \texttt{GĀJIENS: T, (3, 1)} \\
            }
        }
    }
     \\ \\

    \theutterance \stepcounter{utterance}  
    & & & \multicolumn{2}{p{0.3\linewidth}}{
        \cellcolor[rgb]{0.9,0.9,0.9}{
            \makecell[{{p{\linewidth}}}]{
                \texttt{\tiny{[GM$|$GM]}}
                \texttt{GĀJIENS: T, (3, 1)} \\
            }
        }
    }
    & & \\ \\

    \theutterance \stepcounter{utterance}  
    & & & \multicolumn{2}{p{0.3\linewidth}}{
        \cellcolor[rgb]{0.9,0.9,0.9}{
            \makecell[{{p{\linewidth}}}]{
                \texttt{\tiny{[GM$|$GM]}}
                \texttt{Sods: (3,1) nav tukšs, bet satur 'P'.} \\
\\ 
\texttt{Jūs kopā esat uzkrājuši 4 no 16 sodiem.} \\
\texttt{Lūdzu, mēģiniet vēlreiz!} \\
            }
        }
    }
    & & \\ \\

    \theutterance \stepcounter{utterance}  
    & & \multicolumn{4}{p{0.6\linewidth}}{
        \cellcolor[rgb]{0.9,0.9,0.9}{
            \makecell[{{p{\linewidth}}}]{
                \texttt{\tiny{[GM$\rangle$P2]}}
                \texttt{Sods: (3,1) nav tukšs, bet satur 'P'.} \\
\\ 
\texttt{Jūs kopā esat uzkrājuši 4 no 16 sodiem.} \\
\texttt{Lūdzu, mēģiniet vēlreiz!} \\
            }
        }
    }
    & \\ \\

    \theutterance \stepcounter{utterance}  
    & & & \multicolumn{4}{p{0.6\linewidth}}{
        \cellcolor[rgb]{0.9,0.9,0.9}{
            \makecell[{{p{\linewidth}}}]{
                \texttt{\tiny{[GM$\langle$P2]}}
                \texttt{SAKI: Nevaru pārvietot T uz (3,1), jo tur atrodas P. Pirmām jāatbrīvo (3,1). Kur mēs vēlamies novietot P? Varbūt P uz (4,1) un O uz (5,1)? Tad T var iet uz (3,1).vai P uz (4,1)?} \\
            }
        }
    }
     \\ \\

    \theutterance \stepcounter{utterance}  
    & & & \multicolumn{2}{p{0.3\linewidth}}{
        \cellcolor[rgb]{0.9,0.9,0.9}{
            \makecell[{{p{\linewidth}}}]{
                \texttt{\tiny{[GM$|$GM]}}
                \texttt{SAKI: Nevaru pārvietot T uz (3,1), jo tur atrodas P. Pirmām jāatbrīvo (3,1). Kur mēs vēlamies novietot P? Varbūt P uz (4,1) un O uz (5,1)? Tad T var iet uz (3,1).vai P uz (4,1)?} \\
            }
        }
    }
    & & \\ \\

    \theutterance \stepcounter{utterance}  
    & & \multicolumn{4}{p{0.6\linewidth}}{
        \cellcolor[rgb]{0.9,0.9,0.9}{
            \makecell[{{p{\linewidth}}}]{
                \texttt{\tiny{[P1$\langle$GM]}}
                \texttt{Jūsu ziņojums ir nodots otram spēlētājam.} \\
\\ 
\\ 
\texttt{Pašlaik jūs spēlējat kārtu 4 no maksimāli 28.} \\
\texttt{Jūs kopā esat uzkrājuši 4 no 16 sodiem.} \\
\texttt{Otrs spēlētājs nosūtīja šo ziņojumu:} \\
\\ 
\\ 
\texttt{"Nevaru pārvietot T uz (3,1), jo tur atrodas P. Pirmām jāatbrīvo (3,1). Kur mēs vēlamies novietot P? Varbūt P uz (4,1) un O uz (5,1)? Tad T var iet uz (3,1).vai P uz (4,1)?"} \\
\\ 
\\ 
\texttt{Kāda ir jūsu nākamā komanda?} \\
            }
        }
    }
    & \\ \\

    \theutterance \stepcounter{utterance}  
    & \multicolumn{4}{p{0.6\linewidth}}{
        \cellcolor[rgb]{0.9,0.9,0.9}{
            \makecell[{{p{\linewidth}}}]{
                \texttt{\tiny{[P1$\rangle$GM]}}
                \texttt{GĀJIENS: P, (4, 1)} \\
            }
        }
    }
    & & \\ \\

    \theutterance \stepcounter{utterance}  
    & & & \multicolumn{2}{p{0.3\linewidth}}{
        \cellcolor[rgb]{0.9,0.9,0.9}{
            \makecell[{{p{\linewidth}}}]{
                \texttt{\tiny{[GM$|$GM]}}
                \texttt{GĀJIENS: P, (4, 1)} \\
            }
        }
    }
    & & \\ \\

    \theutterance \stepcounter{utterance}  
    & & & \multicolumn{2}{p{0.3\linewidth}}{
        \cellcolor[rgb]{0.9,0.9,0.9}{
            \makecell[{{p{\linewidth}}}]{
                \texttt{\tiny{[GM$|$GM]}}
                \texttt{P moved to (4, 1)} \\
            }
        }
    }
    & & \\ \\

    \theutterance \stepcounter{utterance}  
    & & \multicolumn{4}{p{0.6\linewidth}}{
        \cellcolor[rgb]{0.9,0.9,0.9}{
            \makecell[{{p{\linewidth}}}]{
                \texttt{\tiny{[GM$\rangle$P2]}}
                \texttt{Jūsu ziņojums ir nodots otram spēlētājam.} \\
\\ 
\\ 
\texttt{Pašlaik jūs spēlējat kārtu 5 no maksimāli 28.} \\
\texttt{Jūs kopā esat uzkrājuši 4 no 16 sodiem.} \\
\texttt{Otrs spēlētājs pārvietoja objektu savā režģī.} \\
\\ 
\\ 
\texttt{Kāda ir jūsu nākamā komanda?} \\
            }
        }
    }
    & \\ \\

    \theutterance \stepcounter{utterance}  
    & & & \multicolumn{4}{p{0.6\linewidth}}{
        \cellcolor[rgb]{0.9,0.9,0.9}{
            \makecell[{{p{\linewidth}}}]{
                \texttt{\tiny{[GM$\langle$P2]}}
                \texttt{GĀJIENS: P, (4, 1)} \\
            }
        }
    }
     \\ \\

    \theutterance \stepcounter{utterance}  
    & & & \multicolumn{2}{p{0.3\linewidth}}{
        \cellcolor[rgb]{0.9,0.9,0.9}{
            \makecell[{{p{\linewidth}}}]{
                \texttt{\tiny{[GM$|$GM]}}
                \texttt{GĀJIENS: P, (4, 1)} \\
            }
        }
    }
    & & \\ \\

    \theutterance \stepcounter{utterance}  
    & & & \multicolumn{2}{p{0.3\linewidth}}{
        \cellcolor[rgb]{0.9,0.9,0.9}{
            \makecell[{{p{\linewidth}}}]{
                \texttt{\tiny{[GM$|$GM]}}
                \texttt{P moved to (4, 1)} \\
            }
        }
    }
    & & \\ \\

    \theutterance \stepcounter{utterance}  
    & & \multicolumn{4}{p{0.6\linewidth}}{
        \cellcolor[rgb]{0.9,0.9,0.9}{
            \makecell[{{p{\linewidth}}}]{
                \texttt{\tiny{[P1$\langle$GM]}}
                \texttt{Veiksmīgi pārvietots 'P' uz (4,1). Jūsu atjauninātais režģis izskatās šādi:} \\
\texttt{```} \\
\texttt{ 1234567} \\
\texttt{╔═══════╗} \\
\texttt{║AT◌P◌◌◌║ 1} \\
\texttt{║◌◌◌◌◌◌D║ 2} \\
\texttt{║◌◌◌◌◌┌─╢ 3} \\
\texttt{╟──┐O◌│◌║ 4} \\
\texttt{║◌◌├──┼─╢ 5} \\
\texttt{║◌S│M◌│◌║ 6} \\
\texttt{║◌◌│◌◌│◌║ 7} \\
\texttt{╚══╧══╧═╝} \\
\texttt{```} \\
\texttt{Pašlaik jūs spēlējat kārtu 5 no maksimāli 28.} \\
\texttt{Jūs kopā esat uzkrājuši 4 no 16 sodiem.} \\
\texttt{Otrs spēlētājs pārvietoja objektu savā režģī.} \\
\\ 
\\ 
\texttt{Kāda ir jūsu nākamā komanda?} \\
            }
        }
    }
    & \\ \\

    \theutterance \stepcounter{utterance}  
    & \multicolumn{4}{p{0.6\linewidth}}{
        \cellcolor[rgb]{0.9,0.9,0.9}{
            \makecell[{{p{\linewidth}}}]{
                \texttt{\tiny{[P1$\rangle$GM]}}
                \texttt{GĀJIENS: T, (3, 1)} \\
            }
        }
    }
    & & \\ \\

    \theutterance \stepcounter{utterance}  
    & & & \multicolumn{2}{p{0.3\linewidth}}{
        \cellcolor[rgb]{0.9,0.9,0.9}{
            \makecell[{{p{\linewidth}}}]{
                \texttt{\tiny{[GM$|$GM]}}
                \texttt{GĀJIENS: T, (3, 1)} \\
            }
        }
    }
    & & \\ \\

    \theutterance \stepcounter{utterance}  
    & & & \multicolumn{2}{p{0.3\linewidth}}{
        \cellcolor[rgb]{0.9,0.9,0.9}{
            \makecell[{{p{\linewidth}}}]{
                \texttt{\tiny{[GM$|$GM]}}
                \texttt{T moved to (3, 1)} \\
            }
        }
    }
    & & \\ \\

    \theutterance \stepcounter{utterance}  
    & & \multicolumn{4}{p{0.6\linewidth}}{
        \cellcolor[rgb]{0.9,0.9,0.9}{
            \makecell[{{p{\linewidth}}}]{
                \texttt{\tiny{[GM$\rangle$P2]}}
                \texttt{Veiksmīgi pārvietots 'P' uz (4,1). Jūsu atjauninātais režģis izskatās šādi:} \\
\texttt{```} \\
\texttt{ 1234567} \\
\texttt{╔═══════╗} \\
\texttt{║AD◌P◌◌O║ 1} \\
\texttt{║S◌◌◌M◌◌║ 2} \\
\texttt{║◌◌◌◌◌┌─╢ 3} \\
\texttt{╟──┐◌◌│◌║ 4} \\
\texttt{║◌◌├──┼─╢ 5} \\
\texttt{║◌◌│◌◌│T║ 6} \\
\texttt{║◌◌│◌◌│◌║ 7} \\
\texttt{╚══╧══╧═╝} \\
\texttt{```} \\
\texttt{Pašlaik jūs spēlējat kārtu 6 no maksimāli 28.} \\
\texttt{Jūs kopā esat uzkrājuši 4 no 16 sodiem.} \\
\texttt{Otrs spēlētājs pārvietoja objektu savā režģī.} \\
\\ 
\\ 
\texttt{Kāda ir jūsu nākamā komanda?} \\
            }
        }
    }
    & \\ \\

    \theutterance \stepcounter{utterance}  
    & & & \multicolumn{4}{p{0.6\linewidth}}{
        \cellcolor[rgb]{0.9,0.9,0.9}{
            \makecell[{{p{\linewidth}}}]{
                \texttt{\tiny{[GM$\langle$P2]}}
                \texttt{GĀJIENS: T, (3, 1)} \\
            }
        }
    }
     \\ \\

    \theutterance \stepcounter{utterance}  
    & & & \multicolumn{2}{p{0.3\linewidth}}{
        \cellcolor[rgb]{0.9,0.9,0.9}{
            \makecell[{{p{\linewidth}}}]{
                \texttt{\tiny{[GM$|$GM]}}
                \texttt{GĀJIENS: T, (3, 1)} \\
            }
        }
    }
    & & \\ \\

    \theutterance \stepcounter{utterance}  
    & & & \multicolumn{2}{p{0.3\linewidth}}{
        \cellcolor[rgb]{0.9,0.9,0.9}{
            \makecell[{{p{\linewidth}}}]{
                \texttt{\tiny{[GM$|$GM]}}
                \texttt{T moved to (3, 1)} \\
            }
        }
    }
    & & \\ \\

    \theutterance \stepcounter{utterance}  
    & & \multicolumn{4}{p{0.6\linewidth}}{
        \cellcolor[rgb]{0.9,0.9,0.9}{
            \makecell[{{p{\linewidth}}}]{
                \texttt{\tiny{[P1$\langle$GM]}}
                \texttt{Veiksmīgi pārvietots 'T' uz (3,1). Jūsu atjauninātais režģis izskatās šādi:} \\
\texttt{```} \\
\texttt{ 1234567} \\
\texttt{╔═══════╗} \\
\texttt{║A◌TP◌◌◌║ 1} \\
\texttt{║◌◌◌◌◌◌D║ 2} \\
\texttt{║◌◌◌◌◌┌─╢ 3} \\
\texttt{╟──┐O◌│◌║ 4} \\
\texttt{║◌◌├──┼─╢ 5} \\
\texttt{║◌S│M◌│◌║ 6} \\
\texttt{║◌◌│◌◌│◌║ 7} \\
\texttt{╚══╧══╧═╝} \\
\texttt{```} \\
\texttt{Pašlaik jūs spēlējat kārtu 6 no maksimāli 28.} \\
\texttt{Jūs kopā esat uzkrājuši 4 no 16 sodiem.} \\
\texttt{Otrs spēlētājs pārvietoja objektu savā režģī.} \\
\\ 
\\ 
\texttt{Kāda ir jūsu nākamā komanda?} \\
            }
        }
    }
    & \\ \\

    \theutterance \stepcounter{utterance}  
    & \multicolumn{4}{p{0.6\linewidth}}{
        \cellcolor[rgb]{0.9,0.9,0.9}{
            \makecell[{{p{\linewidth}}}]{
                \texttt{\tiny{[P1$\rangle$GM]}}
                \texttt{GĀJIENS: D, (2, 1)} \\
            }
        }
    }
    & & \\ \\

    \theutterance \stepcounter{utterance}  
    & & & \multicolumn{2}{p{0.3\linewidth}}{
        \cellcolor[rgb]{0.9,0.9,0.9}{
            \makecell[{{p{\linewidth}}}]{
                \texttt{\tiny{[GM$|$GM]}}
                \texttt{GĀJIENS: D, (2, 1)} \\
            }
        }
    }
    & & \\ \\

    \theutterance \stepcounter{utterance}  
    & & & \multicolumn{2}{p{0.3\linewidth}}{
        \cellcolor[rgb]{0.9,0.9,0.9}{
            \makecell[{{p{\linewidth}}}]{
                \texttt{\tiny{[GM$|$GM]}}
                \texttt{D moved to (2, 1)} \\
            }
        }
    }
    & & \\ \\

    \theutterance \stepcounter{utterance}  
    & & \multicolumn{4}{p{0.6\linewidth}}{
        \cellcolor[rgb]{0.9,0.9,0.9}{
            \makecell[{{p{\linewidth}}}]{
                \texttt{\tiny{[GM$\rangle$P2]}}
                \texttt{Veiksmīgi pārvietots 'T' uz (3,1). Jūsu atjauninātais režģis izskatās šādi:} \\
\texttt{```} \\
\texttt{ 1234567} \\
\texttt{╔═══════╗} \\
\texttt{║ADTP◌◌O║ 1} \\
\texttt{║S◌◌◌M◌◌║ 2} \\
\texttt{║◌◌◌◌◌┌─╢ 3} \\
\texttt{╟──┐◌◌│◌║ 4} \\
\texttt{║◌◌├──┼─╢ 5} \\
\texttt{║◌◌│◌◌│◌║ 6} \\
\texttt{║◌◌│◌◌│◌║ 7} \\
\texttt{╚══╧══╧═╝} \\
\texttt{```} \\
\texttt{Pašlaik jūs spēlējat kārtu 7 no maksimāli 28.} \\
\texttt{Jūs kopā esat uzkrājuši 4 no 16 sodiem.} \\
\texttt{Otrs spēlētājs pārvietoja objektu savā režģī.} \\
\\ 
\\ 
\texttt{Kāda ir jūsu nākamā komanda?} \\
            }
        }
    }
    & \\ \\

    \theutterance \stepcounter{utterance}  
    & & & \multicolumn{4}{p{0.6\linewidth}}{
        \cellcolor[rgb]{0.9,0.9,0.9}{
            \makecell[{{p{\linewidth}}}]{
                \texttt{\tiny{[GM$\langle$P2]}}
                \texttt{GĀJIENS: D, (2, 1)} \\
            }
        }
    }
     \\ \\

    \theutterance \stepcounter{utterance}  
    & & & \multicolumn{2}{p{0.3\linewidth}}{
        \cellcolor[rgb]{0.9,0.9,0.9}{
            \makecell[{{p{\linewidth}}}]{
                \texttt{\tiny{[GM$|$GM]}}
                \texttt{GĀJIENS: D, (2, 1)} \\
            }
        }
    }
    & & \\ \\

    \theutterance \stepcounter{utterance}  
    & & & \multicolumn{2}{p{0.3\linewidth}}{
        \cellcolor[rgb]{0.9,0.9,0.9}{
            \makecell[{{p{\linewidth}}}]{
                \texttt{\tiny{[GM$|$GM]}}
                \texttt{Sods: (2,1) nav tukšs, bet satur 'D'.} \\
\\ 
\texttt{Jūs kopā esat uzkrājuši 5 no 16 sodiem.} \\
\texttt{Lūdzu, mēģiniet vēlreiz!} \\
            }
        }
    }
    & & \\ \\

    \theutterance \stepcounter{utterance}  
    & & \multicolumn{4}{p{0.6\linewidth}}{
        \cellcolor[rgb]{0.9,0.9,0.9}{
            \makecell[{{p{\linewidth}}}]{
                \texttt{\tiny{[GM$\rangle$P2]}}
                \texttt{Sods: (2,1) nav tukšs, bet satur 'D'.} \\
\\ 
\texttt{Jūs kopā esat uzkrājuši 5 no 16 sodiem.} \\
\texttt{Lūdzu, mēģiniet vēlreiz!} \\
            }
        }
    }
    & \\ \\

    \theutterance \stepcounter{utterance}  
    & & & \multicolumn{4}{p{0.6\linewidth}}{
        \cellcolor[rgb]{0.9,0.9,0.9}{
            \makecell[{{p{\linewidth}}}]{
                \texttt{\tiny{[GM$\langle$P2]}}
                \texttt{SAKI: Es jau iepriekšējos gājienos pārvietoju D uz (2,1). Redzams, D tagad atrodas (2,1). Mēs jau esam sasniedzuši vienošo stāvokli par A(1,1), D(2,1), T(3,1), P(4,1). Nākamais mērķis: O uz (5,1), S uz (6,1), M uz (7,1). Es nākamajā gājienā pārvietoju O uz (5,1). Tu turpini ar S un M.} \\
            }
        }
    }
     \\ \\

    \theutterance \stepcounter{utterance}  
    & & & \multicolumn{2}{p{0.3\linewidth}}{
        \cellcolor[rgb]{0.9,0.9,0.9}{
            \makecell[{{p{\linewidth}}}]{
                \texttt{\tiny{[GM$|$GM]}}
                \texttt{SAKI: Es jau iepriekšējos gājienos pārvietoju D uz (2,1). Redzams, D tagad atrodas (2,1). Mēs jau esam sasniedzuši vienošo stāvokli par A(1,1), D(2,1), T(3,1), P(4,1). Nākamais mērķis: O uz (5,1), S uz (6,1), M uz (7,1). Es nākamajā gājienā pārvietoju O uz (5,1). Tu turpini ar S un M.} \\
            }
        }
    }
    & & \\ \\

    \theutterance \stepcounter{utterance}  
    & & \multicolumn{4}{p{0.6\linewidth}}{
        \cellcolor[rgb]{0.9,0.9,0.9}{
            \makecell[{{p{\linewidth}}}]{
                \texttt{\tiny{[P1$\langle$GM]}}
                \texttt{Veiksmīgi pārvietots 'D' uz (2,1). Jūsu atjauninātais režģis izskatās šādi:} \\
\texttt{```} \\
\texttt{ 1234567} \\
\texttt{╔═══════╗} \\
\texttt{║ADTP◌◌◌║ 1} \\
\texttt{║◌◌◌◌◌◌◌║ 2} \\
\texttt{║◌◌◌◌◌┌─╢ 3} \\
\texttt{╟──┐O◌│◌║ 4} \\
\texttt{║◌◌├──┼─╢ 5} \\
\texttt{║◌S│M◌│◌║ 6} \\
\texttt{║◌◌│◌◌│◌║ 7} \\
\texttt{╚══╧══╧═╝} \\
\texttt{```} \\
\texttt{Pašlaik jūs spēlējat kārtu 7 no maksimāli 28.} \\
\texttt{Jūs kopā esat uzkrājuši 5 no 16 sodiem.} \\
\texttt{Otrs spēlētājs nosūtīja šo ziņojumu:} \\
\\ 
\\ 
\texttt{"Es jau iepriekšējos gājienos pārvietoju D uz (2,1). Redzams, D tagad atrodas (2,1). Mēs jau esam sasniedzuši vienošo stāvokli par A(1,1), D(2,1), T(3,1), P(4,1). Nākamais mērķis: O uz (5,1), S uz (6,1), M uz (7,1). Es nākamajā gājienā pārvietoju O uz (5,1). Tu turpini ar S un M."} \\
\\ 
\\ 
\texttt{Kāda ir jūsu nākamā komanda?} \\
            }
        }
    }
    & \\ \\

    \theutterance \stepcounter{utterance}  
    & \multicolumn{4}{p{0.6\linewidth}}{
        \cellcolor[rgb]{0.9,0.9,0.9}{
            \makecell[{{p{\linewidth}}}]{
                \texttt{\tiny{[P1$\rangle$GM]}}
                \texttt{GĀJIENS: O, (5, 1)} \\
            }
        }
    }
    & & \\ \\

    \theutterance \stepcounter{utterance}  
    & & & \multicolumn{2}{p{0.3\linewidth}}{
        \cellcolor[rgb]{0.9,0.9,0.9}{
            \makecell[{{p{\linewidth}}}]{
                \texttt{\tiny{[GM$|$GM]}}
                \texttt{GĀJIENS: O, (5, 1)} \\
            }
        }
    }
    & & \\ \\

    \theutterance \stepcounter{utterance}  
    & & & \multicolumn{2}{p{0.3\linewidth}}{
        \cellcolor[rgb]{0.9,0.9,0.9}{
            \makecell[{{p{\linewidth}}}]{
                \texttt{\tiny{[GM$|$GM]}}
                \texttt{O moved to (5, 1)} \\
            }
        }
    }
    & & \\ \\

    \theutterance \stepcounter{utterance}  
    & & \multicolumn{4}{p{0.6\linewidth}}{
        \cellcolor[rgb]{0.9,0.9,0.9}{
            \makecell[{{p{\linewidth}}}]{
                \texttt{\tiny{[GM$\rangle$P2]}}
                \texttt{Jūsu ziņojums ir nodots otram spēlētājam.} \\
\\ 
\\ 
\texttt{Pašlaik jūs spēlējat kārtu 8 no maksimāli 28.} \\
\texttt{Jūs kopā esat uzkrājuši 5 no 16 sodiem.} \\
\texttt{Otrs spēlētājs pārvietoja objektu savā režģī.} \\
\\ 
\\ 
\texttt{Kāda ir jūsu nākamā komanda?} \\
            }
        }
    }
    & \\ \\

    \theutterance \stepcounter{utterance}  
    & & & \multicolumn{4}{p{0.6\linewidth}}{
        \cellcolor[rgb]{0.9,0.9,0.9}{
            \makecell[{{p{\linewidth}}}]{
                \texttt{\tiny{[GM$\langle$P2]}}
                \texttt{GĀJIENS: O, (5, 1)} \\
            }
        }
    }
     \\ \\

    \theutterance \stepcounter{utterance}  
    & & & \multicolumn{2}{p{0.3\linewidth}}{
        \cellcolor[rgb]{0.9,0.9,0.9}{
            \makecell[{{p{\linewidth}}}]{
                \texttt{\tiny{[GM$|$GM]}}
                \texttt{GĀJIENS: O, (5, 1)} \\
            }
        }
    }
    & & \\ \\

    \theutterance \stepcounter{utterance}  
    & & & \multicolumn{2}{p{0.3\linewidth}}{
        \cellcolor[rgb]{0.9,0.9,0.9}{
            \makecell[{{p{\linewidth}}}]{
                \texttt{\tiny{[GM$|$GM]}}
                \texttt{O moved to (5, 1)} \\
            }
        }
    }
    & & \\ \\

    \theutterance \stepcounter{utterance}  
    & & \multicolumn{4}{p{0.6\linewidth}}{
        \cellcolor[rgb]{0.9,0.9,0.9}{
            \makecell[{{p{\linewidth}}}]{
                \texttt{\tiny{[P1$\langle$GM]}}
                \texttt{Veiksmīgi pārvietots 'O' uz (5,1). Jūsu atjauninātais režģis izskatās šādi:} \\
\texttt{```} \\
\texttt{ 1234567} \\
\texttt{╔═══════╗} \\
\texttt{║ADTPO◌◌║ 1} \\
\texttt{║◌◌◌◌◌◌◌║ 2} \\
\texttt{║◌◌◌◌◌┌─╢ 3} \\
\texttt{╟──┐◌◌│◌║ 4} \\
\texttt{║◌◌├──┼─╢ 5} \\
\texttt{║◌S│M◌│◌║ 6} \\
\texttt{║◌◌│◌◌│◌║ 7} \\
\texttt{╚══╧══╧═╝} \\
\texttt{```} \\
\texttt{Pašlaik jūs spēlējat kārtu 8 no maksimāli 28.} \\
\texttt{Jūs kopā esat uzkrājuši 5 no 16 sodiem.} \\
\texttt{Otrs spēlētājs pārvietoja objektu savā režģī.} \\
\\ 
\\ 
\texttt{Kāda ir jūsu nākamā komanda?} \\
            }
        }
    }
    & \\ \\

    \theutterance \stepcounter{utterance}  
    & \multicolumn{4}{p{0.6\linewidth}}{
        \cellcolor[rgb]{0.9,0.9,0.9}{
            \makecell[{{p{\linewidth}}}]{
                \texttt{\tiny{[P1$\rangle$GM]}}
                \texttt{GĀJIENS: S, (6, 1)} \\
            }
        }
    }
    & & \\ \\

    \theutterance \stepcounter{utterance}  
    & & & \multicolumn{2}{p{0.3\linewidth}}{
        \cellcolor[rgb]{0.9,0.9,0.9}{
            \makecell[{{p{\linewidth}}}]{
                \texttt{\tiny{[GM$|$GM]}}
                \texttt{GĀJIENS: S, (6, 1)} \\
            }
        }
    }
    & & \\ \\

    \theutterance \stepcounter{utterance}  
    & & & \multicolumn{2}{p{0.3\linewidth}}{
        \cellcolor[rgb]{0.9,0.9,0.9}{
            \makecell[{{p{\linewidth}}}]{
                \texttt{\tiny{[GM$|$GM]}}
                \texttt{S moved to (6, 1)} \\
            }
        }
    }
    & & \\ \\

    \theutterance \stepcounter{utterance}  
    & & \multicolumn{4}{p{0.6\linewidth}}{
        \cellcolor[rgb]{0.9,0.9,0.9}{
            \makecell[{{p{\linewidth}}}]{
                \texttt{\tiny{[GM$\rangle$P2]}}
                \texttt{Veiksmīgi pārvietots 'O' uz (5,1). Jūsu atjauninātais režģis izskatās šādi:} \\
\texttt{```} \\
\texttt{ 1234567} \\
\texttt{╔═══════╗} \\
\texttt{║ADTPO◌◌║ 1} \\
\texttt{║S◌◌◌M◌◌║ 2} \\
\texttt{║◌◌◌◌◌┌─╢ 3} \\
\texttt{╟──┐◌◌│◌║ 4} \\
\texttt{║◌◌├──┼─╢ 5} \\
\texttt{║◌◌│◌◌│◌║ 6} \\
\texttt{║◌◌│◌◌│◌║ 7} \\
\texttt{╚══╧══╧═╝} \\
\texttt{```} \\
\texttt{Pašlaik jūs spēlējat kārtu 9 no maksimāli 28.} \\
\texttt{Jūs kopā esat uzkrājuši 5 no 16 sodiem.} \\
\texttt{Otrs spēlētājs pārvietoja objektu savā režģī.} \\
\\ 
\\ 
\texttt{Kāda ir jūsu nākamā komanda?} \\
            }
        }
    }
    & \\ \\

    \theutterance \stepcounter{utterance}  
    & & & \multicolumn{4}{p{0.6\linewidth}}{
        \cellcolor[rgb]{0.9,0.9,0.9}{
            \makecell[{{p{\linewidth}}}]{
                \texttt{\tiny{[GM$\langle$P2]}}
                \texttt{GĀJIENS: S, (6, 1)} \\
            }
        }
    }
     \\ \\

    \theutterance \stepcounter{utterance}  
    & & & \multicolumn{2}{p{0.3\linewidth}}{
        \cellcolor[rgb]{0.9,0.9,0.9}{
            \makecell[{{p{\linewidth}}}]{
                \texttt{\tiny{[GM$|$GM]}}
                \texttt{GĀJIENS: S, (6, 1)} \\
            }
        }
    }
    & & \\ \\

    \theutterance \stepcounter{utterance}  
    & & & \multicolumn{2}{p{0.3\linewidth}}{
        \cellcolor[rgb]{0.9,0.9,0.9}{
            \makecell[{{p{\linewidth}}}]{
                \texttt{\tiny{[GM$|$GM]}}
                \texttt{S moved to (6, 1)} \\
            }
        }
    }
    & & \\ \\

    \theutterance \stepcounter{utterance}  
    & & \multicolumn{4}{p{0.6\linewidth}}{
        \cellcolor[rgb]{0.9,0.9,0.9}{
            \makecell[{{p{\linewidth}}}]{
                \texttt{\tiny{[P1$\langle$GM]}}
                \texttt{Veiksmīgi pārvietots 'S' uz (6,1). Jūsu atjauninātais režģis izskatās šādi:} \\
\texttt{```} \\
\texttt{ 1234567} \\
\texttt{╔═══════╗} \\
\texttt{║ADTPOS◌║ 1} \\
\texttt{║◌◌◌◌◌◌◌║ 2} \\
\texttt{║◌◌◌◌◌┌─╢ 3} \\
\texttt{╟──┐◌◌│◌║ 4} \\
\texttt{║◌◌├──┼─╢ 5} \\
\texttt{║◌◌│M◌│◌║ 6} \\
\texttt{║◌◌│◌◌│◌║ 7} \\
\texttt{╚══╧══╧═╝} \\
\texttt{```} \\
\texttt{Pašlaik jūs spēlējat kārtu 9 no maksimāli 28.} \\
\texttt{Jūs kopā esat uzkrājuši 5 no 16 sodiem.} \\
\texttt{Otrs spēlētājs pārvietoja objektu savā režģī.} \\
\\ 
\\ 
\texttt{Kāda ir jūsu nākamā komanda?} \\
            }
        }
    }
    & \\ \\

    \theutterance \stepcounter{utterance}  
    & \multicolumn{4}{p{0.6\linewidth}}{
        \cellcolor[rgb]{0.9,0.9,0.9}{
            \makecell[{{p{\linewidth}}}]{
                \texttt{\tiny{[P1$\rangle$GM]}}
                \texttt{GĀJIENS: M, (7, 1)} \\
            }
        }
    }
    & & \\ \\

    \theutterance \stepcounter{utterance}  
    & & & \multicolumn{2}{p{0.3\linewidth}}{
        \cellcolor[rgb]{0.9,0.9,0.9}{
            \makecell[{{p{\linewidth}}}]{
                \texttt{\tiny{[GM$|$GM]}}
                \texttt{GĀJIENS: M, (7, 1)} \\
            }
        }
    }
    & & \\ \\

    \theutterance \stepcounter{utterance}  
    & & & \multicolumn{2}{p{0.3\linewidth}}{
        \cellcolor[rgb]{0.9,0.9,0.9}{
            \makecell[{{p{\linewidth}}}]{
                \texttt{\tiny{[GM$|$GM]}}
                \texttt{M moved to (7, 1)} \\
            }
        }
    }
    & & \\ \\

    \theutterance \stepcounter{utterance}  
    & & \multicolumn{4}{p{0.6\linewidth}}{
        \cellcolor[rgb]{0.9,0.9,0.9}{
            \makecell[{{p{\linewidth}}}]{
                \texttt{\tiny{[GM$\rangle$P2]}}
                \texttt{Veiksmīgi pārvietots 'S' uz (6,1). Jūsu atjauninātais režģis izskatās šādi:} \\
\texttt{```} \\
\texttt{ 1234567} \\
\texttt{╔═══════╗} \\
\texttt{║ADTPOS◌║ 1} \\
\texttt{║◌◌◌◌M◌◌║ 2} \\
\texttt{║◌◌◌◌◌┌─╢ 3} \\
\texttt{╟──┐◌◌│◌║ 4} \\
\texttt{║◌◌├──┼─╢ 5} \\
\texttt{║◌◌│◌◌│◌║ 6} \\
\texttt{║◌◌│◌◌│◌║ 7} \\
\texttt{╚══╧══╧═╝} \\
\texttt{```} \\
\texttt{Pašlaik jūs spēlējat kārtu 10 no maksimāli 28.} \\
\texttt{Jūs kopā esat uzkrājuši 5 no 16 sodiem.} \\
\texttt{Otrs spēlētājs pārvietoja objektu savā režģī.} \\
\\ 
\\ 
\texttt{Kāda ir jūsu nākamā komanda?} \\
            }
        }
    }
    & \\ \\

    \theutterance \stepcounter{utterance}  
    & & & \multicolumn{4}{p{0.6\linewidth}}{
        \cellcolor[rgb]{0.9,0.9,0.9}{
            \makecell[{{p{\linewidth}}}]{
                \texttt{\tiny{[GM$\langle$P2]}}
                \texttt{GĀJIENS: M, (7, 1)} \\
            }
        }
    }
     \\ \\

    \theutterance \stepcounter{utterance}  
    & & & \multicolumn{2}{p{0.3\linewidth}}{
        \cellcolor[rgb]{0.9,0.9,0.9}{
            \makecell[{{p{\linewidth}}}]{
                \texttt{\tiny{[GM$|$GM]}}
                \texttt{GĀJIENS: M, (7, 1)} \\
            }
        }
    }
    & & \\ \\

    \theutterance \stepcounter{utterance}  
    & & & \multicolumn{2}{p{0.3\linewidth}}{
        \cellcolor[rgb]{0.9,0.9,0.9}{
            \makecell[{{p{\linewidth}}}]{
                \texttt{\tiny{[GM$|$GM]}}
                \texttt{M moved to (7, 1)} \\
            }
        }
    }
    & & \\ \\

    \theutterance \stepcounter{utterance}  
    & & \multicolumn{4}{p{0.6\linewidth}}{
        \cellcolor[rgb]{0.9,0.9,0.9}{
            \makecell[{{p{\linewidth}}}]{
                \texttt{\tiny{[P1$\langle$GM]}}
                \texttt{Veiksmīgi pārvietots 'M' uz (7,1). Jūsu atjauninātais režģis izskatās šādi:} \\
\texttt{```} \\
\texttt{ 1234567} \\
\texttt{╔═══════╗} \\
\texttt{║ADTPOSM║ 1} \\
\texttt{║◌◌◌◌◌◌◌║ 2} \\
\texttt{║◌◌◌◌◌┌─╢ 3} \\
\texttt{╟──┐◌◌│◌║ 4} \\
\texttt{║◌◌├──┼─╢ 5} \\
\texttt{║◌◌│◌◌│◌║ 6} \\
\texttt{║◌◌│◌◌│◌║ 7} \\
\texttt{╚══╧══╧═╝} \\
\texttt{```} \\
\texttt{Pašlaik jūs spēlējat kārtu 10 no maksimāli 28.} \\
\texttt{Jūs kopā esat uzkrājuši 5 no 16 sodiem.} \\
\texttt{Otrs spēlētājs pārvietoja objektu savā režģī.} \\
\\ 
\\ 
\texttt{Kāda ir jūsu nākamā komanda?} \\
            }
        }
    }
    & \\ \\

    \theutterance \stepcounter{utterance}  
    & \multicolumn{4}{p{0.6\linewidth}}{
        \cellcolor[rgb]{0.9,0.9,0.9}{
            \makecell[{{p{\linewidth}}}]{
                \texttt{\tiny{[P1$\rangle$GM]}}
                \texttt{SAKI: pabeigts?} \\
            }
        }
    }
    & & \\ \\

    \theutterance \stepcounter{utterance}  
    & & & \multicolumn{2}{p{0.3\linewidth}}{
        \cellcolor[rgb]{0.9,0.9,0.9}{
            \makecell[{{p{\linewidth}}}]{
                \texttt{\tiny{[GM$|$GM]}}
                \texttt{SAKI: pabeigts?} \\
            }
        }
    }
    & & \\ \\

    \theutterance \stepcounter{utterance}  
    & & \multicolumn{4}{p{0.6\linewidth}}{
        \cellcolor[rgb]{0.9,0.9,0.9}{
            \makecell[{{p{\linewidth}}}]{
                \texttt{\tiny{[GM$\rangle$P2]}}
                \texttt{Veiksmīgi pārvietots 'M' uz (7,1). Jūsu atjauninātais režģis izskatās šādi:} \\
\texttt{```} \\
\texttt{ 1234567} \\
\texttt{╔═══════╗} \\
\texttt{║ADTPOSM║ 1} \\
\texttt{║◌◌◌◌◌◌◌║ 2} \\
\texttt{║◌◌◌◌◌┌─╢ 3} \\
\texttt{╟──┐◌◌│◌║ 4} \\
\texttt{║◌◌├──┼─╢ 5} \\
\texttt{║◌◌│◌◌│◌║ 6} \\
\texttt{║◌◌│◌◌│◌║ 7} \\
\texttt{╚══╧══╧═╝} \\
\texttt{```} \\
\texttt{Pašlaik jūs spēlējat kārtu 11 no maksimāli 28.} \\
\texttt{Jūs kopā esat uzkrājuši 5 no 16 sodiem.} \\
\texttt{Otrs spēlētājs nosūtīja šo ziņojumu:} \\
\\ 
\\ 
\texttt{"pabeigts?"} \\
\\ 
\\ 
\texttt{Kāda ir jūsu nākamā komanda?} \\
            }
        }
    }
    & \\ \\

    \theutterance \stepcounter{utterance}  
    & & & \multicolumn{4}{p{0.6\linewidth}}{
        \cellcolor[rgb]{0.9,0.9,0.9}{
            \makecell[{{p{\linewidth}}}]{
                \texttt{\tiny{[GM$\langle$P2]}}
                \texttt{SAKI: pabeigts!} \\
            }
        }
    }
     \\ \\

    \theutterance \stepcounter{utterance}  
    & & & \multicolumn{2}{p{0.3\linewidth}}{
        \cellcolor[rgb]{0.9,0.9,0.9}{
            \makecell[{{p{\linewidth}}}]{
                \texttt{\tiny{[GM$|$GM]}}
                \texttt{SAKI: pabeigts!} \\
            }
        }
    }
    & & \\ \\

    \theutterance \stepcounter{utterance}  
    & & & \multicolumn{2}{p{0.3\linewidth}}{
        \cellcolor[rgb]{0.9,0.9,0.9}{
            \makecell[{{p{\linewidth}}}]{
                \texttt{\tiny{[GM$|$GM]}}
                \texttt{true} \\
            }
        }
    }
    & & \\ \\

    \theutterance \stepcounter{utterance}  
    & & & \multicolumn{2}{p{0.3\linewidth}}{
        \cellcolor[rgb]{0.9,0.9,0.9}{
            \makecell[{{p{\linewidth}}}]{
                \texttt{\tiny{[GM$|$GM]}}
                \texttt{Initial states:} \\
\texttt{```} \\
\texttt{Player 1:} \\
\texttt{ 1234567} \\
\texttt{╔═══════╗} \\
\texttt{║◌T◌◌◌◌◌║ 1} \\
\texttt{║A◌◌◌◌◌D║ 2} \\
\texttt{║◌◌◌◌◌┌─╢ 3} \\
\texttt{╟──┐O◌│◌║ 4} \\
\texttt{║◌P├──┼─╢ 5} \\
\texttt{║◌S│M◌│◌║ 6} \\
\texttt{║◌◌│◌◌│◌║ 7} \\
\texttt{╚══╧══╧═╝} \\
\\ 
\\ 
\texttt{Player 2:} \\
\texttt{ 1234567} \\
\texttt{╔═══════╗} \\
\texttt{║◌◌P◌A◌O║ 1} \\
\texttt{║S◌◌◌M◌◌║ 2} \\
\texttt{║◌◌◌◌D┌─╢ 3} \\
\texttt{╟──┐◌◌│◌║ 4} \\
\texttt{║◌◌├──┼─╢ 5} \\
\texttt{║◌◌│◌◌│T║ 6} \\
\texttt{║◌◌│◌◌│◌║ 7} \\
\texttt{╚══╧══╧═╝} \\
\\ 
\texttt{```} \\
            }
        }
    }
    & & \\ \\

    \theutterance \stepcounter{utterance}  
    & & & \multicolumn{2}{p{0.3\linewidth}}{
        \cellcolor[rgb]{0.9,0.9,0.9}{
            \makecell[{{p{\linewidth}}}]{
                \texttt{\tiny{[GM$|$GM]}}
                \texttt{End states:} \\
\texttt{```} \\
\texttt{Player 1:} \\
\texttt{ 1234567} \\
\texttt{╔═══════╗} \\
\texttt{║ADTPOSM║ 1} \\
\texttt{║◌◌◌◌◌◌◌║ 2} \\
\texttt{║◌◌◌◌◌┌─╢ 3} \\
\texttt{╟──┐◌◌│◌║ 4} \\
\texttt{║◌◌├──┼─╢ 5} \\
\texttt{║◌◌│◌◌│◌║ 6} \\
\texttt{║◌◌│◌◌│◌║ 7} \\
\texttt{╚══╧══╧═╝} \\
\\ 
\\ 
\texttt{Player 2:} \\
\texttt{ 1234567} \\
\texttt{╔═══════╗} \\
\texttt{║ADTPOSM║ 1} \\
\texttt{║◌◌◌◌◌◌◌║ 2} \\
\texttt{║◌◌◌◌◌┌─╢ 3} \\
\texttt{╟──┐◌◌│◌║ 4} \\
\texttt{║◌◌├──┼─╢ 5} \\
\texttt{║◌◌│◌◌│◌║ 6} \\
\texttt{║◌◌│◌◌│◌║ 7} \\
\texttt{╚══╧══╧═╝} \\
\\ 
\texttt{```} \\
            }
        }
    }
    & & \\ \\

    \theutterance \stepcounter{utterance}  
    & & & \multicolumn{2}{p{0.3\linewidth}}{
        \cellcolor[rgb]{0.9,0.9,0.9}{
            \makecell[{{p{\linewidth}}}]{
                \texttt{\tiny{[GM$|$GM]}}
                \texttt{* success: True} \\
\texttt{* lose: False} \\
\texttt{* aborted: False} \\
\texttt{{-}{-}{-}{-}{-}{-}{-}} \\
\texttt{* Shifts: 7.00} \\
\texttt{* Max Shifts: 12.00} \\
\texttt{* Min Shifts: 6.00} \\
\texttt{* End Distance Sum: 0.00} \\
\texttt{* Init Distance Sum: 30.04} \\
\texttt{* Expected Distance Sum: 29.33} \\
\texttt{* Penalties: 5.00} \\
\texttt{* Max Penalties: 16.00} \\
\texttt{* Rounds: 11.00} \\
\texttt{* Max Rounds: 28.00} \\
\texttt{* Object Count: 7.00} \\
            }
        }
    }
    & & \\ \\

\end{supertabular}
}

\end{document}
